% GENERATED FILE. Edit canonical lessons, then run npm run generate:books.
% canonical-source-hash: fnv1a64:88c35bf46589760e
% canonical-lessons: SA-C139-bhavisyat, SA-C139-svah, SA-C139-agami, SA-C139-bhavisyati, SA-C139-gamisyati

\chapter{The Future, Tomorrow, Coming, Will Be, Will Go}
\label{ch:sa-the-future-tomorrow-coming-will-be-will-go}
\hlchaptermodality{139}

\begin{chapteropening}
\textbf{By the end of this chapter:} \emph{I can say what will happen, and when: tomorrow, next week, next year.}

Complete the last lesson of chapter 139: I can say what will happen, and when: tomorrow, next week, next year.
\end{chapteropening}

\section[\texorpdfstring{\sk{भविष्यत्} (bhaviṣyat)}{भविष्यत् (bhaviṣyat)}]{\sk{भविष्यत्} (bhaviṣyat) — the future}

\label{lesson:SA-C139-bhavisyat}



\begin{quote}
Before the new one: say the Sanskrit for wanders, then the Sanskrit for grows.
\end{quote}

\subsection*{You'll want to know: \sk{भविष्यत्}}
\textbf{\sk{भविष्यत्}} — \emph{bhaviṣyat} — \textquotedblleft{}the future\textquotedblright{}. Sanskrit puts the future inside the verb, with \textbf{-\sk{ष्य}-} before the ending: \textbf{\sk{पठति}} — \textquotedblleft{}reads\textquotedblright{} — becomes \textbf{\sk{पठिष्यति}} — \textquotedblleft{}will read\textquotedblright{}.

\begin{grammarlens}[title={What you've built}]
One more word for this chapter.
\end{grammarlens}

\subsection*{Your turn}
\begin{itemize}
\raggedright
  \item \emph{Say it:} \emph{bhaviṣyat}
  \item \emph{Say it:} \emph{bhaviṣyat}, once more
  \item \emph{Say it:} say \emph{bhaviṣyat}
  \item \emph{Recall:} say the Sanskrit for wanders, then the Sanskrit for grows, then say \emph{bhaviṣyat} again
\end{itemize}

\begin{tcolorbox}[breakable,colback=teal!4,colframe=teal!35!black,title={Before you move on}]
What does \sk{भविष्यत्} mean? (\textquotedblleft{}The future\textquotedblright{}.) Say it once more.
\end{tcolorbox}
\section[\texorpdfstring{\sk{श्वः} (śvaḥ)}{श्वः (śvaḥ)}]{\sk{श्वः} (śvaḥ) — tomorrow}

\label{lesson:SA-C139-svah}



\begin{quote}
Before the new one: say the Sanskrit for grows, then the Sanskrit for the future.
\end{quote}

\subsection*{You'll want to know: \sk{श्वः}}
\textbf{\sk{श्वः}} — \emph{śvaḥ} — \textquotedblleft{}tomorrow\textquotedblright{}. The time word stands first and the future follows: \textbf{\sk{श्वः} \sk{सः} \sk{पठिष्यति।}} — \textquotedblleft{}Tomorrow he will read.\textquotedblright{} You already have \textbf{\sk{परश्वः}}, the day after it.

\begin{grammarlens}[title={What you've built}]
One more word for this chapter.
\end{grammarlens}

\subsection*{Your turn}
\begin{itemize}
\raggedright
  \item \emph{Say it:} \emph{śvaḥ}
  \item \emph{Say it:} \emph{śvaḥ}, once more
  \item \emph{Say it:} say \emph{śvaḥ}
  \item \emph{Recall:} say the Sanskrit for grows, then the Sanskrit for the future, then say \emph{śvaḥ} again
\end{itemize}

\begin{tcolorbox}[breakable,colback=teal!4,colframe=teal!35!black,title={Before you move on}]
What does \sk{श्वः} mean? (\textquotedblleft{}Tomorrow\textquotedblright{}.) Say it once more.
\end{tcolorbox}
\section[\texorpdfstring{\sk{आगामी} (āgāmī)}{आगामी (āgāmī)}]{\sk{आगामी} (āgāmī) — coming, next}

\label{lesson:SA-C139-agami}



\begin{quote}
Before the new one: say the Sanskrit for the future, then the Sanskrit for tomorrow.
\end{quote}

\subsection*{You'll want to know: \sk{आगामी}}
\textbf{\sk{आगामी}} — \emph{āgāmī} — \textquotedblleft{}coming, next\textquotedblright{}. It stands before the time word and takes its ending: \textbf{\sk{आगामिनि} \sk{सप्ताहे}} — \textquotedblleft{}next week\textquotedblright{}; \textbf{\sk{आगामिनि} \sk{संवत्सरे}} — \textquotedblleft{}next year\textquotedblright{}.

\begin{grammarlens}[title={What you've built}]
One more word for this chapter.
\end{grammarlens}

\subsection*{Your turn}
\begin{itemize}
\raggedright
  \item \emph{Say it:} \emph{āgāmī}
  \item \emph{Say it:} \emph{āgāmī}, once more
  \item \emph{Say it:} say \emph{āgāmī}
  \item \emph{Recall:} say the Sanskrit for the future, then the Sanskrit for tomorrow, then say \emph{āgāmī} again
\end{itemize}

\begin{tcolorbox}[breakable,colback=teal!4,colframe=teal!35!black,title={Before you move on}]
What does \sk{आगामी} mean? (\textquotedblleft{}Coming, next\textquotedblright{}.) Say it once more.
\end{tcolorbox}
\section[\texorpdfstring{\sk{भविष्यति} (bhaviṣyati)}{भविष्यति (bhaviṣyati)}]{\sk{भविष्यति} (bhaviṣyati) — will be}

\label{lesson:SA-C139-bhavisyati}



\begin{quote}
Before the new one: say the Sanskrit for tomorrow, then the Sanskrit for coming.
\end{quote}

\subsection*{You'll want to know: \sk{भविष्यति}}
\textbf{\sk{भविष्यति}} — \emph{bhaviṣyati} — \textquotedblleft{}will be\textquotedblright{}. The future of \textbf{\sk{भवति}}, \textquotedblleft{}is, becomes\textquotedblright{}: \textbf{\sk{श्वः} \sk{रविवासरः} \sk{भविष्यति।}} — \textquotedblleft{}Tomorrow will be Sunday.\textquotedblright{}

\begin{grammarlens}[title={What you've built}]
One more word for this chapter.
\end{grammarlens}

\subsection*{Your turn}
\begin{itemize}
\raggedright
  \item \emph{Say it:} \emph{bhaviṣyati}
  \item \emph{Say it:} \emph{bhaviṣyati}, once more
  \item \emph{Say it:} say \emph{bhaviṣyati}
  \item \emph{Recall:} say the Sanskrit for tomorrow, then the Sanskrit for coming, then say \emph{bhaviṣyati} again
\end{itemize}

\begin{tcolorbox}[breakable,colback=teal!4,colframe=teal!35!black,title={Before you move on}]
What does \sk{भविष्यति} mean? (\textquotedblleft{}Will be\textquotedblright{}.) Say it once more.
\end{tcolorbox}
\section[\texorpdfstring{\sk{गमिष्यति} (gamiṣyati)}{गमिष्यति (gamiṣyati)}]{\sk{गमिष्यति} (gamiṣyati) — will go}

\label{lesson:SA-C139-gamisyati}



\begin{quote}
Before the new one: say the Sanskrit for coming, then the Sanskrit for will be.
\end{quote}

\subsection*{You'll want to know: \sk{गमिष्यति}}
\textbf{\sk{गमिष्यति}} — \emph{gamiṣyati} — \textquotedblleft{}will go\textquotedblright{}. You met \textbf{\sk{अगच्छत्}}, \textquotedblleft{}he went\textquotedblright{}; this is the same verb looking forward: \textbf{\sk{परश्वः} \sk{सः} \sk{आपणं} \sk{गमिष्यति।}} — \textquotedblleft{}The day after tomorrow he will go to the shop.\textquotedblright{}

\begin{grammarlens}[title={What you've built}]
One more word for this chapter.
\end{grammarlens}

\subsection*{Your turn}
\begin{itemize}
\raggedright
  \item \emph{Say it:} \emph{gamiṣyati}
  \item \emph{Say it:} \emph{gamiṣyati}, once more
  \item \emph{Say it:} say \emph{gamiṣyati}
  \item \emph{Recall:} say the Sanskrit for coming, then the Sanskrit for will be, then say \emph{gamiṣyati} again
\end{itemize}

\begin{tcolorbox}[breakable,colback=teal!4,colframe=teal!35!black,title={Before you move on}]
What does \sk{गमिष्यति} mean? (\textquotedblleft{}Will go\textquotedblright{}.) Say it once more.
\end{tcolorbox}
